\begin{progprob}
    Commercial flights must balance the weight of passengers between the two
    sides of the plane in order to ensure a safe takeoff. Suppose you are piloting a
    full flight and must choose a passenger from the left and a passenger from
    the right to switch places so that the total weight of the passengers on
    the left side is as close as possible to the total weight of the passengers
    on the right.

    In a file named \python{balance.py}, design an algorithm named \python{balance(left, right)}
    which takes in two \emph{sorted} arrays, \python{left} and \python{right}, containing the weights
    of the passengers on the left side of the plane and the right side of the plane, respectively.
    Your function should return a pair \python{(left_ix, right_ix)} containing
    the index of the person on  the left and the index of the person on the
    right who should be swapped to bring the weights to as close as balanced as
    possible. If the left and right sides are already balanced, return
    \python{None}. If there are multiple optimal solutions (that is, multiple
    ways of swapping passengers that result in the same optimal balance),
    return just one of them.

    For example, suppose \python{left = [4, 8, 12]} and \python{right = [6, 15]}.
    When \python{balance(left, right)} is called, your function should return
    \python{(1, 0)} because the passenger with weight 8 on the left should swap
    with the passenger with weight 6 on the right. This would result in the
    left side having a total weight of 4 + 6 + 12 = 24 and the right side
    having a total weight of 8 + 15 = 23, which is as close as possible.

    To receive full credit, your algorithm must be optimal in the sense that its worst
    case time complexity is as good as theoretically possible.

    \begin{soln}

        \inputminted{python}{\thisdir/include-solution/balance.py}

    \end{soln}
\end{progprob}
